% Section 1.5, from lectures/L01/README.md: the objectives, the evaluation questions, and the next
% lecture.

\needspace{10\baselineskip}
\section{Review}
\label{c1:sec:review}

\subsection*{What you should be able to do now}
After this chapter, you should:
\begin{itemize}
  \item Know how \code{std::vector} is used in C++.
  \item Know why a traditional program isn't well suited for mimicking human intelligence.
  \item Know how machine learning differs from traditional programming in terms of input, output,
    and rules.
  \item Be able to train simple linear regression models by hand.
  \item Be able to implement training of a linear regression model in software.
\end{itemize}

\needspace{6\baselineskip}
\subsection*{Questions to test yourself}
You should be able to explain:
\begin{enumerate}
  \item What's the difference between input, rules, and output in a traditional program compared
    to a machine learning system?
  \item What's meant by an epoch when training a machine learning model?
  \item How are the parameters in a linear regression model used to compute a prediction for a
    given input?
  \item Why is an interface used for the linear regression model? What does it enable?
  \item What happens to the model's weight and bias value when input \code{x} equals zero during
    training?
  \item How does the learning rate affect training? What happens if it's set too high or too low?
\end{enumerate}

\needspace{6\baselineskip}
\subsection*{Looking ahead}
Chapter~\ref{c2:ch:training} continues the implementation of regression models in software:
\begin{itemize}
  \item Randomizing the training order.
  \item Computing the precision of a trained model.
\end{itemize}
